% ===== 原文 pp.82--86 =====
% [原文 p.82]
由耦合项 (4.27) 可见，作用量被修改为
\begin{equation}
S\to S+2\pi s(d)\int \Lambda^m_{d-3}\cup_\eta\delta w^m_2\,,\qquad s(d)=(-1)^d\,,
\tag{4.112}
\end{equation}
而利用 (4.107) 可知，这蕴含着电与磁对称性之间的混合 't Hooft 反常
\begin{equation}
Z\left[B_2^e,B^m_{d-2}+d\Lambda^m_{d-3}\right]=\exp\left(-2\pi\ii s(d)\int\Lambda^m_{d-3}\cup_\eta\mathrm{Bock}(B_2^e)\right)\times Z\left[B_2^e,B^m_{d-2}\right].
\tag{4.113}
\end{equation}
于是反常理论为
\begin{equation}
\mathcal{I}_{d+1}[B_2^e,B^m_{d-2}]=\exp\left(2\pi\ii s(d)\int B^m_{d-2}\cup_\eta\mathrm{Bock}(B_2^e)\right).
\tag{4.114}
\end{equation}

% [Tutor Note: p.82 式 (4.112)——原文此处把磁障碍类写成 $w^m_2$（带上标 $m$），而 (4.107)、(4.115)、(4.135) 以及例 4.8 中的 (4.127) 都写成 $w_2$；按上下文二者是同一个 $\mathcal{Z}$-值 2-上链（障碍类），本译本照原文抄录，并用批注提示这一记号不一致。判据：若 $w^m_2$ 真是另一个上链，则 (4.112) 与 (4.107) 联立就得不出 (4.113) 中的 $\mathrm{Bock}(B_2^e)$。]

\begin{notebox}{Bockstein 的消失：分裂的短正合列}
在反常的表达式 (4.114) 中，只有 $\mathrm{Bock}(B_2^e)$ 的群上同调类才是相关的。每当该上同调类消失
\begin{equation}
[\mathrm{Bock}(B_2^e)]=0\,,
\tag{4.115}
\end{equation}
反常也就消失。这发生在短正合列 (4.100) \textbf{分裂}时，即若能找到同态
\begin{equation}
s:\ G^{(1)}\to\mathcal{X}\,,
\tag{4.116}
\end{equation}
使得
\begin{equation}
\pi\circ s=1\,.
\tag{4.117}
\end{equation}
存在满足上述性质的同态 $s$，等价于把 $\mathcal{X}$ 分解为
\begin{equation}
\mathcal{X}=\mathcal{Z}\times\mathcal{Y}\,,
\tag{4.118}
\end{equation}
其中 $\mathcal{Y}\cong G^{(1)}$。
\end{notebox}

\begin{examplebox}{例 4.8：纯 $\SO(6)$ Yang--Mills 理论}
作为一个简单例子，考虑规范群为
\begin{equation}
\mathcal{G}=\SO(6)\,,
\tag{4.119}
\end{equation}
且没有物质内容的规范理论。单连通群为
\begin{equation}
G=\Spin(6)=\SU(4)\,,
\tag{4.120}
\end{equation}
其中心为
\begin{equation}
Z(G)=\Z_4\,,
\tag{4.121}
\end{equation}
并且我们有
\begin{equation}
\mathcal{Z}=\Z_2\,,
\tag{4.122}
\end{equation}
% [原文 p.83]
因为
\begin{equation}
\SO(6)=\Spin(6)/\Z_2\,.
\tag{4.123}
\end{equation}
磁 $(d-3)$-形式对称性为
\begin{equation}
G^{(d-3)}=\hat{\mathcal{Z}}=\Z_2\,,
\tag{4.124}
\end{equation}
而电 1-形式对称性为
\begin{equation}
G^{(1)}=Z(\mathcal{G})=Z(G)/\mathcal{Z}=\Z_4/\Z_2=\Z_2\,.
\tag{4.125}
\end{equation}
这些群落在短正合列
\begin{equation}
\begin{aligned}
0&\to\mathcal{Z}\to\mathcal{X}=Z(G)\to G^{(1)}\to 0\\
&=0\to\Z_2\to\Z_4\to\Z_2\to 0\,.
\end{aligned}
\tag{4.126}
\end{equation}
开启一个 $\Z_2$-值磁性背景场 $B^m_{d-2}$ 会引入耦合项
\begin{equation}
\pi\int B^m_{d-2}\cup w_2\,,
\tag{4.127}
\end{equation}
其中 $w_2$ 是把 $\SO(6)$ 丛提升为 $\Spin(6)$ 丛的 $\Z_2$-值障碍，$\cup$ 是标准的 cup 积。另一方面，$\Z_2$-值背景场 $B_2^e$ 是把
\begin{equation}
\PSO(6)=\SO(6)/\Z_2
\tag{4.128}
\end{equation}
丛提升为 $\SO(6)$ 丛的障碍。总体而言，把 $\P\SO(6)$ 丛提升为 $\Spin(6)$ 丛有一个 $\Z_4$-值障碍 $W_2$，它可表示为
\begin{equation}
W_2=2w_2+\tilde{B}_2^e\,,
\tag{4.129}
\end{equation}
其中单射
\begin{equation}
i:\ \Z_2\to\Z_4
\tag{4.130}
\end{equation}
就是乘以 2，即
\begin{equation}
\begin{aligned}
2w_2(v_0,v_1,v_2)&=0 &&\text{若 } w_2(v_0,v_1,v_2)=0\,,\\
2w_2(v_0,v_1,v_2)&=2 &&\text{若 } w_2(v_0,v_1,v_2)=1\,,
\end{aligned}
\tag{4.131}
\end{equation}
以及
\begin{equation}
\begin{aligned}
\tilde{B}_2^e(v_0,v_1,v_2)&\in\{0,2\} &&\text{若 } B_2^e(v_0,v_1,v_2)=0\,,\\
\tilde{B}_2^e(v_0,v_1,v_2)&\in\{1,3\} &&\text{若 } B_2^e(v_0,v_1,v_2)=1\,.
\end{aligned}
\tag{4.132}
\end{equation}
可以验证
\begin{equation}
\delta\tilde{B}_2^e(v_0,v_1,v_2,v_3)\in\{0,2\}\,,
\tag{4.133}
\end{equation}
从而得到
\begin{equation}
\delta\tilde{B}_2^e\equiv 2\mathrm{Bock}(B_2^e)\,,
\tag{4.134}
\end{equation}
% [原文 p.84]
其中 $\mathrm{Bock}(B_2^e)$ 是一个 $\Z_2$-值 3-上链。障碍类 $w_2$ 是非闭的，
\begin{equation}
\delta w_2=-\mathrm{Bock}(B_2^e)\,,
\tag{4.135}
\end{equation}
而反常理论 (4.114) 为
\begin{equation}
\mathcal{I}_{d+1}[B_2^e,B^m_{d-2}]=\exp\left(\pi\ii\int B^m_{d-2}\cup\mathrm{Bock}(B_2^e)\right),
\tag{4.136}
\end{equation}
因为对 $\Z_2$-值量而言符号 $s(d)$ 是无关紧要的。
\end{examplebox}

% [Tutor Note: p.83 式 (4.126)--(4.134)——这一段是 $\Z_2\to\Z_4\to\Z_2$ 型 Bockstein 的具体计算，可用来核对自己的 Bockstein 公式：$i$ 是乘 2 的单射，故 $\delta\tilde{B}_2^e=2\,\mathrm{Bock}(B_2^e)$ 必为偶值（式 (4.134)），这正是 (4.133) 的内容。锚点：若把 $\tilde{B}_2^e$ 的取值误写成 $\{0,1\}$，则 (4.133) 会得到 $1$，与 (4.134) 矛盾。]

% [Tutor Note: p.84 式 (4.136)——注意到这里 $2\pi\ii\,s(d)\to\pi\ii$：因为 $B^m_{d-2}\cup\mathrm{Bock}(B_2^e)$ 是 $\Z_2$-值且 $s(d)=\pm1$，$(2\pi s(d))\bmod 2\pi$ 只依赖于 $s(d)$ 的符号，而对 $\Z_2$-值量 $\exp(\pm\pi\ii\,n)=\exp(\pi\ii\,n)$。所以原文说「$s(d)$ 无关紧要」。]

\subsection{规范化（gauging）}

\subsubsection{定义}

规范化（gauging）是这样一种程序：把 $p$-形式对称性的背景场 $B_{p+1}$ 提升为动力学的 $(p+1)$-形式（或 $(p+1)$-上链）规范场 $b_{p+1}$
\begin{equation}
B_{p+1}\to b_{p+1}\,.
\tag{4.137}
\end{equation}
具体地说，这意味着我们要对所有可能的规范不等价背景 $B_{p+1}$ 求和。也就是说，我们在背景场的每个等价类（模背景规范变换）中选取一个代表 $B_{p+1}$，并对所有这样的代表求和。若原来的理论是 $\mathcal{T}$，而被规范化的 $p$-形式对称群是 $G^{(p)}$，则规范化后得到的理论称为\textbf{规范化理论}（gauged theory），记作
\begin{equation}
\mathcal{T}/G^{(p)}\,.
\tag{4.138}
\end{equation}
规范化理论的配分函数为
\begin{equation}
\begin{aligned}
Z_{\mathcal{T}/G^{(p)}}&\propto\sum_{[B_{p+1}]}Z_{\mathcal{T}}[B_{p+1}],\qquad G^{(p)}\ \text{离散},\\[2pt]
Z_{\mathcal{T}/G^{(p)}}&\propto\int\mathcal{D}[B_{p+1}]\,Z_{\mathcal{T}}[B_{p+1}],\qquad G^{(p)}\ \text{连续}.
\end{aligned}
\tag{4.139}
\end{equation}
上述程序只有在代表选取无关紧要时才是合理的，即要求
\begin{equation}
Z_{\mathcal{T}}[B^g_{p+1}]=Z_{\mathcal{T}}[B_{p+1}]\,,
\tag{4.140}
\end{equation}
其中 $B^g_{p+1}$ 与 $B_{p+1}$ 由如 (4.35) 的规范变换相联系；这导致如下陈述。

\begin{statementbox}{陈述 4.4：'t Hooft 反常是规范化的障碍}
一个 $G^{(p)}$ $p$-形式对称群只有在它不带 't Hooft 反常时才能被规范化。
\end{statementbox}

% [Tutor Note: p.84 式 (4.139)——离散情形求和下标是 $[B_{p+1}]$（规范等价类），连续情形是 $\mathcal{D}[B_{p+1}]$（泛函积分测度）；这两者的区别正是「离散群的对偶变量是离散的、连续群的对偶是连续的」。核对锚点：$G^{(p)}$ 离散时式 (4.139) 与 (4.158) 的求和形式必须一致（后者只是在每项前乘一个系数 $I_d[B_{p+1}]$），否则离散挠率那一节就对不上。]

% [原文 p.85]
\para{非阿贝尔规范理论中电 1-形式对称性的规范化。}利用上述求和的规范化描述，很容易理解非阿贝尔规范理论中电 1-形式对称性的规范化。我们先看一个例子，然后给出一般结果。

\begin{examplebox}{例 4.9：从 $\SU(2)$ 到 $\SO(3)$}
考虑 $d$ 维纯 $\SU(2)$ Yang--Mills 理论，它具有电 1-形式对称性
\begin{equation}
G^{(1)}=\Z_2\,.
\tag{4.141}
\end{equation}
正如例 4.3 中讨论的，在存在背景场 $B_2^e$ 时，$\SU(2)$ 规范理论要对 $\SO(3)$ 丛的一个子集求和，即那些提升为 $\SU(2)$ 丛的障碍为
\begin{equation}
w_2=B_2^e\,.
\tag{4.142}
\end{equation}
的丛。对 $\Z_2$ 1-形式对称性作规范化后，我们把 $w_2$ 提升为动力学场，并对所有 $\SO(3)$ 丛求和，而不论其障碍为何。于是规范化后的理论是 $d$ 维纯 $\SO(3)$ Yang--Mills 理论
\begin{equation}
\mathcal{T}_{\SU(2)}/\Z_2^{(1)}=\mathcal{T}_{\SO(3)}\,.
\tag{4.143}
\end{equation}
\end{examplebox}

我们可以容易地推广上述论证。考虑一个 $d$ 维非阿贝尔规范理论 $\mathcal{T}_{\mathcal{G}}$，其规范群为 $\mathcal{G}$，并带有某些物质内容。电 1-形式对称群为
\begin{equation}
G^{(1)}\subseteq Z(\mathcal{G})\,.
\tag{4.144}
\end{equation}
考虑规范化它的一个子群
\begin{equation}
H^{(1)}\subseteq G^{(1)}\,.
\tag{4.145}
\end{equation}
在存在 $H^{(1)}$ 1-形式对称性的背景场 $B_2^e$ 时，$\mathcal{G}$ 规范理论要对那些 $B_2^e$ 恰为提升为 $\mathcal{G}$ 丛之障碍的 $\mathcal{G}/H$ 丛求和。对所有 $B_2^e$ 求和意味着我们现在对所有 $\mathcal{G}/H$ 丛求和，因此规范化 $H^{(1)}$ 后得到的理论 $\mathcal{T}_{\mathcal{G}}/H^{(1)}$ 与 $\mathcal{T}_{\mathcal{G}}$ 的唯一差别在于 $\mathcal{T}_{\mathcal{G}}/H^{(1)}$ 的规范群是
\begin{equation}
\mathcal{G}/H^{(1)}\,.
\tag{4.146}
\end{equation}
我们可以把得到的理论记为 $\mathcal{T}_{\mathcal{G}/H^{(1)}}$。于是规范化的结果可表示为
\begin{equation}
\mathcal{T}_{\mathcal{G}}/H^{(1)}=\mathcal{T}_{\mathcal{G}/H^{(1)}}\,.
\tag{4.147}
\end{equation}

% [Tutor Note: p.85 式 (4.143)、(4.147)——注意两种上下标写法的区别：$\mathcal{T}_{\SU(2)}/\Z_2^{(1)}$ 表示「对 $\mathcal{T}_{\SU(2)}$ 规范化 $\Z_2^{(1)}$」，而 $\mathcal{T}_{\mathcal{G}/H^{(1)}}$ 表示「规范群为 $\mathcal{G}/H^{(1)}$ 的理论」；(4.147) 声明二者相等。核对锚点：取 $\mathcal{G}=\SU(2)$、$H^{(1)}=\Z_2^{(1)}$，则 (4.147) 退化为 (4.143)：$\SU(2)/\Z_2=\SO(3)$，故两式的对应关系自洽。]

% [原文 p.86]
\para{非阿贝尔规范理论中磁 $(d-3)$-形式对称性的规范化。}类似地，也很容易理解非阿贝尔规范理论中磁 $(d-3)$-形式对称性的规范化。我们同样先看一个例子，然后给出一般结果。

\begin{examplebox}{例 4.10：从 $\SO(3)$ 到 $\SU(2)$}
考虑 $d$ 维纯 $\SO(3)$ Yang--Mills 理论，它具有磁 $(d-3)$-形式对称性
\begin{equation}
G^{(d-3)}=\Z_2\,.
\tag{4.148}
\end{equation}
正如例 4.4 中讨论的，在存在背景场 $B^m_{d-2}$ 时，$\SO(3)$ 规范理论的作用量包含一项
\begin{equation}
\pi\int B^m_{d-2}\cup w_2\,,
\tag{4.149}
\end{equation}
其中 $w_2$ 是动力学场。对 $B^m_{d-2}$ 求和时，上述项充当 Lagrange 乘子，它强制
\begin{equation}
w_2=0\,.
\tag{4.150}
\end{equation}
换言之，所得的路径积分只对 $\SU(2)$ 规范丛求和。因此，规范化后的理论是 $d$ 维纯 $\SU(2)$ Yang--Mills 理论
\begin{equation}
\mathcal{T}_{\SO(3)}/\Z_2^{(d-3)}=\mathcal{T}_{\SU(2)}\,.
\tag{4.151}
\end{equation}
\end{examplebox}

我们可以容易地推广上述论证。考虑一个 $d$ 维非阿贝尔规范理论 $\mathcal{T}_{\mathcal{G}}$，其规范群为 $\mathcal{G}$，并带有某些物质内容。磁 $(d-3)$-形式对称性为
\begin{equation}
G^{(d-3)}=\hat{\mathcal{Z}}\,,
\tag{4.152}
\end{equation}
其中 $\mathcal{Z}$ 通过
\begin{equation}
\mathcal{G}=G/\mathcal{Z}
\tag{4.153}
\end{equation}
把 $\mathcal{G}$ 与单连通群 $G$ 联系起来。考虑规范化它的一个子群
\begin{equation}
H^{(d-3)}\subseteq G^{(d-3)}\,.
\tag{4.154}
\end{equation}
运用 Lagrange 乘子论证可知，这一规范化程序把规范群改为
\begin{equation}
G/\mathcal{Y}\,,
\tag{4.155}
\end{equation}
其中
\begin{equation}
\mathcal{Y}=\widehat{G^{(d-3)}/H^{(d-3)}}\subseteq\mathcal{Z}\,.
\tag{4.156}
\end{equation}
于是我们有
\begin{equation}
\mathcal{T}_{\mathcal{G}}/H^{(d-3)}=\mathcal{T}_{G/\mathcal{Y}}\,.
\tag{4.157}
\end{equation}

% [Tutor Note: p.86 式 (4.156)——$\mathcal{Y}$ 定义为 $G^{(d-3)}/H^{(d-3)}$ 的 Pontryagin 对偶（$\widehat{\ \cdot\ }$ 表示对偶群 $\mathrm{Hom}(\cdot,\U(1))$），并断言它落在 $\mathcal{Z}$ 中。可核对的特例：取 $G^{(d-3)}=\Z_4$、$H^{(d-3)}=\Z_2$，则商为 $\Z_2$，其对偶仍为 $\Z_2$，且 $\Z_2\subseteq\Z_4$ 成立。注意对有限 阿贝尔群总有 $\widehat{G}\cong G$（$\Z_n$ 自对偶，直积保持这一性质），所以本例中 $\widehat{\Z_2}=\Z_2$ 是合法的。英文原词：Pontryagin dual group（Pontryagin 对偶群）。]
